\chapter{Gerätelehre --- \gAasRsN{RS~IX}}

Thomas Hochreiter [R?]; BEMAW, Sebastian Gabriel

Derzeit in Vorbereitung. 

Kapitelliste nicht vollst\"andig

\section{Blutdruckmessung}

\label{topic:rr-messung}

Grundlagen

Fallstricke

Die Blutdruckmanschette soll nach M\"oglichkeit nicht auf der Seite

\begin{itemize}
    \item eines Shuntarmes bei Dialysepatienten
    \begin{itemize}
        \item Ein Dialyse-Shunt ist ein k\"unstlicher Kurzschluss zwischen einer
        Armarterie und einer Armvene. Dadurch wird die Vene erweitert und sie
        pulsiert, mann kann viel mehr Durchfluss f\"ur die Dialyse erreichen.
        Die Nahtstelle zwischen Vene und Arterie ist jedoch sehr sensibel, wenn
        gestaut wird kann sie aufplatzen 
        \item ${\rightarrow}$ Gefahr einer massiven Shuntblutung
    \end{itemize}
    \item eines Arm\"odems
    \item einer Brust-OP mit Lymphknoten-Entfernung (\"Odemgefahr)
    \item einer Halbseitenschw\"ache oder -l\"ahmung
\end{itemize}

angelegt werden. 

\gCave{Shunt-Arme sind tabu f\"ur die Blutdruckmessung!}

Bitten Sie im Praktikum am KTW Dialysepatienten h\"oflich Ihnen ihren
Shunt-Arm zu zeigen!

Ma{\ss}nahmen

\begin{enumerate}
    \item Evt. Funktion des Blutdruckmessger\"ates und des Stethoskops pr\"ufen
    \begin{itemize}
        \item Kontrolle der Dichte des Ventils, der Verbindungsschl\"auche und
        der Manschette (Ventil schlie{\ss}en, Manschette etwas aufpumpen und
        Zusammenpressen, Manschette ganz entl\"uften)
        \item Kontrolle der G\"angigkeit des Ventils (es soll sich leicht auf
        und zuschrauben lassen)
        \item Kontrolle des Manometers auf korrekte Eichung: Die Nadel muss sich
        im Ruhezustand genau bei 0\,mmHg befinden
        \item Ohrst\"opsel des Stethoskops zeigen schr\"ag nach vorne
        \item Trichterfunktion durch Ber\"uhren der Membran pr\"ufen
    \end{itemize}
    \item Patientenlagerung
    \begin{itemize}
        \item Arm freimachen und entspannt auf einen Tisch legen in leichter
        Flexion und Supination
        \item Oberarm auf Herzh\"ohe
    \end{itemize}
    \item Auswahl der geeigneten Manschettengr\"o{\ss}e
    \begin{itemize}
        \item Manschette laut Herstellerangaben w\"ahlen z.B Standardmanschette:
        Oberarmumfang max. 31 cm
        \item Messwertverf\"alschung vermeiden! Der aufblasbare Teil der
        Manschette soll 80\% des Oberarmumfangs umfassen, die Breite der
        Manschette mindestens 1/3 des Oberarmumfangs betragen
        \item Zu kleine Manschetten f\"uhren zum Messen zu hoher Werte, zu weite
        Manschetten k\"onnen nicht am Arm fixiert werden
    \end{itemize}
    \item Handhabung des Stethoskops
    \item Bestimmen der Anlegestelle f\"ur Manschette und der
    Auskultationsstelle durch Palpation
    \begin{itemize}
        \item Arteria brachialis medial in der Ellenbeuge aufsuchen
    \end{itemize}
    
    \item Anlegen der Manschette
    \begin{itemize}
        \item Manschette muss auf unbekleideter Haut aufliegen
        \item m\"a{\ss}ig straff
        \item Abstand zur Ellenbeuge 2,5 cm
        \item aufblasbaren Teil \"uber A. brachialis zentrieren
        \item Schl\"auche nicht im Weg,
        \item CAVE: Halbseitenschw\"ache/-l\"ahmung, Shuntarm bei
        Dialysepatienten, Arm\"odem
    \end{itemize}
    \item Puls f\"uhlen beim Aufpumpen
    \begin{itemize}
        \item Schnelles Aufpumpen
        \item Tasten des Radialispulses mit Finger 2-4
        \item Aufpumpen bis 30 mmHg \"uber jenen Druck, ab dem kein Radialispuls
        mehr getastet wird
    \end{itemize}
    \item Messvorgang
    \begin{itemize}
        \item Kopf des Stethoskops auf a. brachialis legen
        \item Druck langsam ablassen 2 mmHg/sec
        \item Systolischer Wert: Beginn der turbulenten Korotkoffschen
        Str\"omungsger\"ausche
        \item diastolischer Wert: Verschwinden der Korotkoffschen Ger\"ausche
        \item Wenn Ger\"ausche \"uberhaupt nicht verschwinden (z.B. bei Kindern,
        Schwangeren etc.), so gilt als diastolischer Wert der Messpunkt, an dem
        die Ger\"ausche deutlich leiser werden
        \item Wenn Ger\"ausche verschwunden sind, noch weitere 10mmHg langsam
        ablassen -- wenn dann keine Ger\"ausche mehr h\"orbar sind, kann die
        Manschette z\"ugig entleert werden
        \item CAVE: auskultatorische L\"ucke
    \end{itemize}
    \item Protokollierung der Blutdruckwerte
\end{enumerate}


\section{Das EKG}

Das Elektrokardiogramm (EKG) zeigt mittels eines entsprechenden
Ger\"ates die elektrische Herzaktivit\"at des
Reizbildungs- und Reizleitungssystems an. Dazu werden auf
den K\"orper nach einem bestimmten Muster Elektroden
aufgeklebt und mittels eines Kabels mit dem EKG-Ger\"at
verbunden. Im Rettungsdienst werden fast ausschließlich
Klebeelektroden für den einmaligen Gebrauch verwendet. 

\gAuthNotes{GABS: Sollte im Defi-Teil durchgenommen
werden?:} \gZ{EKG -- Ablauf einer Erregung

\begin{itemize}
    \item P-Welle \~{} Vorhoferregung
    \item PQ-Zeit \~{} AV-\"Uberleitung
    \item QRS-Komplex \~{} Kammererregung
    \item ST-Strecke + T-Welle \~{} Erregungsr\"uckbildung
\end{itemize}}

\gCave{Das EKG gibt nur die elektrische Leitung, nicht aber
die tats\"achliche Muskelarbeit wieder! \gEs{Es erlaubt
keine Aussage \"uber die Herzleistung! }}

\gFazit{Grundsatz: \gEs{Selbst ein toter Patient kann ein
{\quotedblbase}sch\"ones{\textquotedblleft} EKG haben.}
(allerdings nicht lange)}

\subsection{EKG -- Ableitungen}
Grunds\"atzlich unterscheidet man zwischen den Extremit\"atenableitungen
und den Brustwandableitungen. 

\begin{itemize}
    \item Extremit\"atenableitungen: I, II, III, aVL, aVR, aVF
    \item Brustwandableitungen: V1 -- V6
\end{itemize}

Es haben sich folgende Begriffe eingebürgert:

\begin{description}
    \item[Extramitätenableitung] "`\gT{Rhythmusstreifen}"'
    zur Beurteilung des Herzrhythmus und zum Monitoring
    \item[Extremitätenableitung\,+\,Brustwandableitung]
    "`\gT{12-Kanal-EKG}"' -- Damit können Schäden der
    jeweiligen Herzregion zugeordnet werden. Wenn vorhanden
    ist das 12-Kanal-EKG (\emph{"`12er"'}) für die
    Diagnostik das Mittel der Wahl. 
\end{description}

\subsubsection{Extremit\"atenableitungen}

Die Extremit\"atenableitungen werden mittels 4 Elektroden abgeleitet. Je
eine Elektrobe wird an den Enden der Extremit\"aten angebracht. Zur
Vereinfachung k\"onnen die Elektroden auch am \"Ubergang vom Rumpf zu
den Extremit\"aten angebracht werden (z.B. statt dem rechten Arm an der
rechten Schulter). Wichtig ist das die die
Elektrodenpäarchen immer auf \gEs{gleicher Höhe} angebracht
werden!

Aus den Informationen der 4 Elektroden erzeugt das
EKG-Ger\"at \gZ{6} Ableitungen\gZ{:  I, II, III, aVL, aVR
und aVF.}

\begin{center}
\includegraphics[width=7cm]{./images/ECGcolor.png}
\captionof{figure}{\gAuthNotesPicLegalOk{}{PD}Extremitätenableitungen
(schwarze Kabel) und Brustwandableitungen (blaue Kabel).
\gSourcePicture{Madhero88 (Wikimedia Commons), Public
domain}\gAuthNotes{GABS: Bessere Beschriftung wäre
wünschenswert, Dringlichkeit: niedrig}}
\end{center}


 
\gAuthNotes{Elektroden noch farblich Kennzeichnen}
\begin{center}
\begin{tabulary}{\linewidth}{CCC}\toprule
\gTabHeading{Elektrode} & \gTabHeading{Position} &
\gTabHeading{Alternative Postion}
\\\midrule
\gTabSub{\color{red}Rot} & Rechte Hand & Rechte
Schulter
\\
\gTabSub{\color{yellow}Gelb} & Linke Hand & Linke
Schulter
\\
\gTabSub{\color{darkGreen}Gr\"un } & Linker Fu{\ss}  & Linke
Leiste
\\
\gTabSub{\color{black}Schwarz} & Rechter Fu{\ss} & Rechte
Leiste \\\bottomrule
\end{tabulary}
\end{center}

\subsubsection{Brustwandableitungen}

Die Brustwandableitungen werden mittels 6 Elektroden abgeleitet und
erlauben eine genauere Zuordnung von St\"orungen zu der jeweiligen
Region des Herzens. Aus den Informationen der 6 Elektroden erzeugt das
EKG-Ger\"at 6 Ableitungen: V1 -- V6. 

\begin{center}
\begin{tabulary}{\linewidth}[H]{cL}
\toprule
Elektrode
 &
Position
\\\midrule
V1 & 4. Zwischenrippenraum rechts vom Sternum
\\
V2 & 4. Zwischenrippenraum links vom Sternum
\\
V3 & Mitte zwischen V2 und V4
\\
V4 & 5. Zwischenrippenraum links in der Verl\"angerung der
Schl\"usselbein-Mittellinie
\\
V5 & Mitte zwischen V4 und V6
\\
V6 & 6. Zwischenrippenraum in der mittleren Achsellinie
\\\bottomrule
\end{tabulary}
\end{center}





\section{$O_2$ --
Sauerstoff}\label{topic:sauerstoff}\label{topic:o2}\index{Sauerstoff}\index{$O_2$}

\gDefinition{Sauerstoff (chem. Symbol $O_2$) ist ein
lebenswichtiges Gas welches in unserer Atemluft vorkommt und au{\ss}erdem ein essentielles Notfall{\textquotedblright}medikament{\textquotedblleft}, die
fr\"uhzeitige Verabreichung ist wichtig f\"ur die Prognose des
Patienten!}

Sauerstoff darf und muss bei Bedarf vom RS verabreicht werden.

\begin{center}
\includegraphics[width=4.522cm,height=5.628cm]{./images/rippenspreitzer-02-fett.png}
\captionof{figure}{NEIN! \gSourcePicture{Rippenspreitzer.de}}

\end{center}
\subsection{Sauerstoff -- technische Grundlagen}
Sauerstoff wird in Druckflaschen  gelagert. Laut EU-Norm sind diese
Flaschen wei{\ss} und mit einem {\quotedblbase}N{\textquotedblleft}
gekennzeichnet. Das {\quotedblbase}N{\textquotedblleft} steht f\"ur
{\quotedblbase}neue Kennzeichnung{\textquotedblleft}, da fr\"uher die
Flaschen blau gekennzeichnet wurden. 

Wichtige Begriffe

\begin{itemize}
    \item F\"ullgr\"o{\ss}e
    \item F\"ullvolumen
    \item Durchflu{\ss}
\end{itemize}

Die Druckflaschen haben unterschiedliche F\"ullgr\"o{\ss}en (Volumina),
g\"angige Flaschengr\"o{\ss}en sind 0,5~l, 1~l, 1,5~l, 2~l, 5~l, 10~l.
Der Sauerstoff wird unter Druck gelagert, durch diesen Trick kann die
Menge des zur Verf\"ugung stehenden Sauerstoffs um ein Vielfaches
gesteigert werden. Der Druck einer gef\"ullten Flasche betr\"agt
normalerweise ca. 200\,bar. 

Jede Druckflasche verf\"ugt \"uber ein \gT{Hauptventil} und
ein Standardgewinde. An diesem Gewinde ist entweder ein
Druckschlauch oder ein \gT{Druckminderer} angeschlossen. Der 
Druckminderer hat ein \gE{Abgabeventil}, welches den
\gE{Druck drosselt} und eine Abgabe an den Patienten
erm\"oglicht. Je nach Bauart kann der Druckminderer \"uber
ein \gT{Regelventil}, ein \gT{Flowmeter} und ein
\gT{Barometer} verf\"ugen. Mit dem \gE{Regelventil} kann man
die Durchlu{\ss}- bzw. Abgaberate ein\-stellen, das
\gE{Flowmeter} (Durchflu{\ss}anzeige) zeigt die aktuelle
Durchflu{\ss}rate in Liter pro Minute (\gLMin) an. Das
\gE{Barometer} zeigt den Druck in der Sauerstofflasche an
wenn das Haupt\-ventil ge\"offnet ist.

\begin{center}
\includegraphics[width=7.189cm,height=4.94cm]{./images/beatmungsgeraet-medumat.png}
\captionof{figure}{Beatmungsgerät \emph{Medumat
Standard\textregistered}. \gSourcePicture{ASB
Floridsdorf-Donaustadt}}
\end{center}

Das effektive F\"ullvolumen erh\"alt man wenn man die F\"ullgr\"o{\ss}e
(in~l) mit dem Flaschendruck (in bar) multipliziert.  

Wenn man wissen m\"ochte wie lange die vorhandene Sauerstoffmenge bei
einer bestimmten Abgabemenge reicht dividiert man einfach die
vorhandene Menge durch die Abgaberate. Achtung: Es ist eine
Ausreichende Reserve einzuplanen, und es ist zu Bedenken dass in der
Flasche ein Restdruck verbleiben muss! 

\subsection{Heiteres Rechnen mit Sauerstoff}

Mit Sauerstoff l\"asst sich gut rechnen ...

\begin{equation}
\mbox{Verfügbare Menge in l} = {\mbox{Volumen der Flasche}}
\times {\mbox{Druck in bar}}
\end{equation}

oder kurz: 

\begin{equation}
l = {Volumen} \times {Druck}
\end{equation}



\minisec{Beispiel}

\begin{quotation}
Flasche 10\,l mit 150\,bar: 1500\,l  $O_2$

Patient ben\"otigt 5\,\unitfrac{l}{min} $O_2$ : 

Die F\"ullung reicht f\"ur 300 Minuten, abz\"uglich eines
Sicherheitsvolumens.
\end{quotation}

\subsection{Umgang mit Sauerstoff und Druckflaschen}

Die Sauerstoffflaschen sind pfleglich zu behandeln

Sauerstoff ist zusammen mit Fett selbstentz\"undlich! Die Armaturen und
Gewinde sind fettfrei zu halten! Das Hantieren mit Feuer und offenem
Licht ist in der N\"ahe von Sauerstoffflaschen verboten! 

Der gro{\ss}e Druck birgt weitere Gefahren: Sollte ein Ventil abbrechen
kann die Flasche eine enorme Kraft entwickeln und zu einer Rakete
werden. Die Kraft einer mit 200\,bar gef\"ullten Flasche reicht aus um
Mauern zu druchbrechen!  

Achtung! Vorsicht beim Hantieren -  Explosionsgefahr!

Flaschen stehen unter Druck, ein Abbrechen der Ventile ist katastrophal!

Sauerstoff ist zusammen mit Fett selbstentz\"undlich ${\rightarrow}$
Ventile nicht mit fetten/\"oligen H\"anden bedienen.

\subsection{Verabreichungsarten}
\subsubsection{Sauerstoffbrille}
Nur bei vorhandener Spontanatmung!

zwei kurze, in die Nasen\"offnungen  eingef\"uhrte Kunst\-stoff\-stutzen

bei ca. 4\,l O2 /min bis 40\% $O_2$  in Atemgas  m\"oglich


Bei behinderter Nasenatmung kann die  Sauerstoffbrille im Mundwinkel
plaziert werden. 

\begin{center}
\begin{tabulary}{\linewidth}[H]{CLLL}
\hline
Was
 &
Wann
 &
Vorteile
&
Nachteile
\\\hline 
\gT{Sauerstoff\-brille}
 &
Spontanatmung

$O_2$-Flow {\textless} 5\gLMin
 &
gute Toleranz

fehlendes Engegef\"uhl

Sprechen\&Husten m\"oglich
 &
ungenaue Dosierung

Austrocknung  der Schleimh\"aute
\\\hline
\gT{Sauerstoff\-maske}
 &
Spontanatmung

$O_2$-Flow ab 5\gLMin

 &
h\"ohere $O_2$-Konzentration
 &
h\"ohere $O_2$-Konzentration

$CO_2$-R\"uckatmung  wenn $O_2$-Durchflu{\ss}
{\textless} 5\gLMin!) \\\hline
\gT{Sauerstoff\-maske mit Reservoir}
 &
 &
Noch h\"ohere $O_2$-Konzentration
 &
Wie oben.
\\\hline
\gT{Beatmungs\-beutel} mit  Reservoir
 &
\gEs{Beatmung} des Patienten 
 &
 &
\\\hline
\end{tabulary}
\end{center}

\gDias
{\includegraphics[width=2.5cm]{./images/o2-brille.png}\captionof{figure}{Sauerstoffbrille}}
{\includegraphics[width=2.5cm]{./images/o2-maske.png}\captionof{figure}{Sauerstoffmaske}}
{\includegraphics[width=2.5cm]{./images/o2-maske-reservoir.png}
\captionof{figure}{Sauerstoffmaske mit Reservoir}}

\subsubsection{Beatmung mit Beatmungsbeutel}

\gCave{JEDE BEATMUNG SOLLTE MIT RESERVOIRBEUTEL UND 10-15\gLMin
$O_2$ DURCHGEF\"UHRT WERDEN!}

Beatmungsfrequenz beim Erwachsenen: 12-15\,x\,/\,min (Eigenfrequenz)

Achtung! Zuviel ist schlecht! Nicht zu schnell beatmen ${\rightarrow}$
Hyperventilationsgefahr!

\subsection{Sauerstofftherapie}
Prinzipiell sollte sich die gew\"ahlte Therapie an der Diagnose und am
Schweregrad orientieren!

Bedenke: 

\begin{itemize}
\item Diagnose!

\item Schockbek\"ampfung

\item Isch\"amie (bei Akutem Koronoarsyndrom, Insult, ..)

\item Atemst\"orungen (Fremdk\"orper, Atemnot, {\dots})

\end{itemize}

\subsubsection{Pulsoxymetrie}

\index{Pulsoxymetrie}\begin{center}
 [Warning: Image ignored] % Unhandled or unsupported graphics:
%\includegraphics[width=3.516cm,height=3.516cm]{AASdev20090228-img13.png}
\gTempPicMissing{Pulsoxy img13}

\end{center}
Das Pulsoxymeter (kurz: Pulsoxy) mi{\ss}t dden Puls und die
Sauerstoffs\"attigung (\index{SpO2}SpO2)des Blutes (Verh\"altnis
O2-beladene/ nichtbeladene Erythrozyten). 

Normalwert 96-99\%, 

[F078?] Achtung: Ein {\quotedblbase}guter{\textquotedblleft}
\index{SpO2}Sp$O_2$-Wert bedeutet nicht {\quotedblbase}keinen Sauerstoff
geben{\textquotedblleft}. Die Entscheidung bez\"uglich Sauerstoffgabe
ber\"ucksichtigt die Diagnose und den Zustand des Patienten, der
\index{SpO2}Sp$O_2$-Wert hat wenig bis gar nichts zu sagen! 

Achtung: Das Pulsoxymeter verwechselt $O_2$ und
$CO$-Molek\"ule, bei einer $CO$-Vergiftung zeigt das Pulsoxy
falsch-hohe S\"attigungswerte an!

Weitere Fehlerquellen: kein Signal: Nagellack, Zentralisation

Das Pulsoxymeter misst die Herzfrequenz und die Sauerstoffs\"attigung
(=Sp$O_2$) des Blutes, aber es gilt:

[F078?] Ein {\quotedblbase}guter{\textquotedblleft} S\"attigungswert ist
keine Begr\"undung keinen Sauerstoff zu verabreichen!
